% Front-end tutorial fragment. Source excerpts refer to the Rustless checkout.
\chapter{同一个程序，为什么需要不同表示}
\label{chap:front-representations}

\section*{导读}
阅读程序时，我们会同时关注它的写法、含义和执行过程。编译器也需要这些视角，但它不能像读者一样依靠经验补足缺失的信息：一个名字指向什么声明，一次运算使用什么类型，一条分支通往哪里，都必须有明确的表示。本章从一个简单的计算开始，说明编译器为什么逐步改变程序的表示，以及每次改变应当解决什么问题。后续章节中的 AST、Visitor 和语义类型，都可以放回这条主线来理解。

\section{从一个计算开始}
考虑下面的 Rustless 函数。Rustless 是本教程讨论的语言子集；它采用若干 Rust 风格的语法，但不具有完整 Rust 的所有语言能力。
\begin{lstlisting}[language=Rustless]
fn calculate(a: i32, b: i32, c: i32) -> i32 {
    a + b * c
}
\end{lstlisting}
源文本按字符从左到右排列。理解这个函数时，我们却不会把这种排列直接当作计算顺序：乘法属于加法的右操作数，整个表达式又是函数体的结果。文本顺序、结构上的包含关系和计算上的依赖关系，已经是三种不同的关系。

如果只保存字符串 \texttt{a + b * c}，每一个使用者都需要重新识别运算符、判断优先级并找到操作数。若编译器首先把乘法与加法的父子关系建立起来，打印器、类型检查器和代码生成器就可以共享这份解释。这是引入程序表示的第一个理由：\emph{已经确定的关系应当成为可以读取的数据，而不应反复隐藏在文本中等待重新解释。}

我们暂时将运算理解为一张依赖图。乘法的结果必须存在，加法才能使用它；参数的名字如何找到声明，则是另一项尚未展开的任务。这样分开观察，可以避免在介绍 AST 时过早讨论寄存器，也避免在讨论机器指令时重新处理括号。

下图呈现前述文字中的计算依赖：树的上方是使用结果的运算，下方是为它提供输入的结构。边表示操作数关系，不表示字符相邻。
\begin{figure}[htbp]
\centering
\begin{tikzpicture}[>=Latex,level distance=13mm,
  every node/.style={draw,rounded corners,font=\small,minimum width=14mm},
  level 1/.style={sibling distance=39mm},level 2/.style={sibling distance=22mm}]
\node {加法}
  child {node {\texttt{a}}}
  child {node {乘法} child {node {\texttt{b}}} child {node {\texttt{c}}}};
\end{tikzpicture}
\caption{表达式的结构明确了计算依赖。}
\label{fig:front-expression-dependency}
\end{figure}
注意，图中的每一条边都有含义。若把两个操作数仅保存为一个没有说明的数组，读者还需要知道第零个元素和第一个元素分别扮演什么角色。稍后采用的 \texttt{left} 与 \texttt{right} 字段，会把这种含义直接写入类型。

\section{词法、语法与语义}
词法分析首先确定程序由哪些词法单元组成。在上例中，\texttt{fn} 是关键字，\texttt{calculate} 是标识符，\texttt{->} 是标点组合，\texttt{i32} 的拼写可以作为路径中的名字出现。词法单元带有类别和位置，因此下一阶段不必逐个字符重新辨认这些单元。

语法分析确定这些单元如何组成语言允许的结构。它能够判断参数列表是否闭合、函数体是否由大括号界定，以及乘法如何嵌入加法。然而，满足语法规则还没有说明名字是否存在，也没有说明两个操作数是否支持同一个运算。
\begin{lstlisting}[language=Rustless]
fn main() {
    let value = 1 + true;
}
\end{lstlisting}
这段代码包含合法的整数、布尔值和加法写法，能够形成一个二元表达式结构。它的问题出现在语义层面：加法需要适合该运算的操作数。类似地，\texttt{unknown + 1} 可以具有清楚的语法结构，但仍需要名字解析来判断 \texttt{unknown} 指向什么。

这里形成了三个逐步增加的判断：字符能否形成词法单元，单元能否形成语法结构，结构能否获得合法的语言含义。将它们分开，使诊断能够指出错误所属的层次，也使某一层的修改不必连带改变所有层。例如，增加一种整数进制主要影响字面量识别与解释；它不应迫使二元表达式改变左右子节点的结构。

真实前端的入口定义位于 \path{E:/RComplier/Rustless/include/rx/frontend.hpp}。下列摘录保留原有接口；阅读时关注返回结果，而不必立即了解 ANTLR 的类。
\begin{lstlisting}[language=C++]
struct FrontendResult {
  std::unique_ptr<const ast::Crate> crate;
  // False means the generated parser rejected the token stream.
  bool syntax_ok{};
};

FrontendResult parse_source(std::string_view source,
                            DiagnosticEngine &diagnostics);
\end{lstlisting}
\texttt{syntax\_ok} 说明解析是否成功；\texttt{crate} 是已建立的语法树。该结果没有把语法成功表述为整个程序已经通过类型检查。后续阶段仍有独立的任务和独立的输入要求。

\section{AST、语义事实、IR 与机器代码}
抽象语法树，即 AST，按语言概念组织程序。一个函数有签名和函数体，一个二元表达式有运算符和两个操作数，一个条件表达式有条件与分支。AST 的首要目的，是让这些语法关系成为明确的结构。

语义分析在这份结构上确定含义。例如，上例中三个参数的类型都是 \texttt{i32}；\texttt{a} 这一次出现与参数声明之间存在引用关系；乘法和加法得到相应的运算规则。语义事实可以与 AST 分开保存，因为同一棵树的结构没有因类型判断而改变。把事实按节点身份组织起来，也便于下一个阶段直接查询已知的结论。

中间表示，即 IR，进一步显式表达计算、数据流、控制流与内存操作。表达式树中隐含的依赖，在 IR 中通常体现为值被哪条指令定义、又被哪些指令使用。以下是解释依赖的教学记法，并非从当前 Rustless 完整 lowering 得到的输出：
\begin{lstlisting}
t1 = multiply(b, c)
t2 = add(a, t1)
return t2
\end{lstlisting}
一个语法节点不必对应恰好一条 IR 指令。结构体构造可能涉及存储分配与多个字段写入；短路逻辑可能形成多个基本块；方法调用还可能需要接收者转换。机器表示随后继续引入目标平台的寄存器、栈槽和调用约定。因此，要求所有阶段共享一套通用节点，往往会把不同阶段的关系混杂在一起。

下图把同一个函数放入不同表示，箭头说明每次转换新增的明确关系。虚线表示当前 Rustless 尚未完成的完整源程序转换与机器后端。
\begin{figure}[htbp]
\centering
\begin{tikzpicture}[>=Latex,node distance=9mm,
  box/.style={draw,rounded corners,text width=95mm,align=center,font=\small}]
\node[box] (text) {源文本：名字、标点和运算的写法};
\node[box,below=of text] (ast) {AST：声明、表达式与结构性子节点};
\node[box,below=of ast] (sem) {语义事实：声明身份、类型与已解析的操作};
\node[box,below=of sem] (ir) {IR：值定义、控制流、内存与调用};
\node[box,below=of ir] (machine) {机器表示：寄存器、栈帧与目标指令};
\draw[->] (text)--node[right,font=\small]{解析与建树}(ast);
\draw[->] (ast)--node[right,font=\small]{语义分析}(sem);
\draw[->,dashed] (sem)--node[right,font=\small]{完整 lowering}(ir);
\draw[->,dashed] (ir)--node[right,font=\small]{后端转换}(machine);
\end{tikzpicture}
\caption{逐步改变表示，使不同关系在适合它们的阶段显式出现。}
\end{figure}
当前项目已经具有 AST、语义事实以及有类型的 IR 模型、构造与验证接口。完整 lowering 和机器后端仍属于后续设计范围。阅读它们时，应当区分已有的数据结构与尚待完成的转换算法。

\section{阶段之间需要传递什么}
当语义分析指出右操作数类型错误时，它需要找到相应的源文本；当 lowering 查询一个表达式时，它需要知道正在查询哪一次出现；当调用已经解析时，它需要知道目标是哪一个声明。这三种需求分别引出源码范围、节点身份和声明身份。

源码范围回答“这一结构来自哪里”。节点身份回答“这一条事实属于哪一次语法出现”。声明身份回答“这个名字或调用最终指向哪个声明”。同样的拼写可以出现在不同范围中，也可以因作用域不同而指向不同声明。若把三者都用一个字符串代替，后续阶段便会丢失必要的区分。

类型与转换也需要分开。一段表达式可以本来具有引用类型，却在某个操作数位置被隐式解引用。语义阶段应当记录该使用位置所需要的转换，后端按记录实施，而不必重新猜测语言规则。源码来源则可以继续伴随生成指令，以服务诊断与观察；来源身份与 IR 值身份仍是不同的概念。

下图追踪一个加法表达式的源码范围、语法身份、语义结论与生成值，说明它们彼此关联，却各自回答不同的问题。
\begin{figure}[htbp]
\centering
\begin{tikzpicture}[>=Latex,node distance=12mm,
  box/.style={draw,rounded corners,align=center,font=\small,text width=37mm}]
\node[box] (n) {AST 节点\\\texttt{NodeId} 与 \texttt{Span}};
\node[box,right=of n] (s) {语义事实\\类型、操作与输入转换};
\node[box,below=of s] (v) {IR 值\\由指令定义的结果};
\node[box,below=of n] (src) {源码切片\\\texttt{a + b * c}};
\draw[->] (n)--node[above,font=\small]{按身份查询}(s);
\draw[->] (n)--node[left,font=\small]{按范围切片}(src);
\draw[->,dashed] (s)--node[right,font=\small]{转换}(v);
\draw[->,dashed] (n)--node[below,font=\small]{来源关系}(v);
\end{tikzpicture}
\caption{一项计算中的不同身份与信息关系。}
\end{figure}
这些信息不能只留在产生它们的局部函数中。清楚的阶段接口，应当使消费者得到完成任务所需的信息，并且能够说明哪些关系已经检查、哪些仍需要处理。后续的具体载荷和语义结果设计，都是这项原则的应用。

\section{用什么评价一种设计}
评价一种表示，可以从结构、决定的归属、扩展方向和代价四方面展开。结构方面，询问字段是否直接表达语言概念：\texttt{left} 与 \texttt{right} 是否比没有角色说明的数组更容易阅读。决定的归属方面，询问同一规则是否在多个阶段重复实现：后端是否还在按字符串猜测操作数类型。

扩展有不同方向。增加一种表达式，需要补充它的数据结构及相关处理规则；增加一种分析过程，则需要解释已有节点的新用途。这两种变化不应混为同一个“可扩展性”问题。第七章讨论 Visitor 时，将用节点与过程的矩阵具体说明这一区别。

运行代价同样需要具体观察。逐节点分配会产生分配和指针访问成本；较大的联合会增大某些节点的对象尺寸；模板可以提供静态检查，也可能增加实例化和编译时间。清楚的数据结构使这些成本更容易分析，但不能替代实际测量。教程中的“优雅”因此不是对一种技术的赞美，而是对关系是否明确、责任是否一致、变化是否可追踪的判断。

本章建立了表示之间的主线。接下来先处理所有表示都可能需要的基础信息：源文本如何保存，以及一个结构怎样准确地指回文本中的位置。

\chapter{源文本、词法单元与诊断位置}
\label{chap:front-source}

\section*{导读}
编译器不仅需要理解程序，还需要把理解结果解释给编写程序的人。错误位置、运算符位置以及保留的字面量拼写，都依赖可靠的源码表示。本章先讨论文本的寿命，再说明词法类别与非法输入，最后联系字符坐标、字节范围和生成代码。学习本章之后，读者应能够解释一个 AST 节点的范围来自哪里，以及为什么词法识别不能包办字面量的全部含义。

\section{SourceFile 与源码寿命}
一个源码切片只有在原文本仍然存在时才有意义。当前 \texttt{Crate} 拥有 \texttt{SourceFile}，而 \texttt{SourceFile} 拥有一个 \texttt{std::string}。AST 的范围因而可以在语法分析器对象销毁之后继续用于诊断与打印。

下面是 \path{E:/RComplier/Rustless/include/rx/ast/item.hpp} 的真实摘录：
\begin{lstlisting}[language=C++]
struct Crate : Node {
  explicit Crate(std::string source) : source(std::move(source)) {}
  SourceFile source;
  std::vector<Item> items;
};
\end{lstlisting}
这里的拥有关系比“AST 临时借用 parser 输入缓冲区”更容易说明。读者可以沿着 \texttt{Crate} 找到文本所有者；使用者只要维持整棵树的寿命，便可以读取其源码。\texttt{SourceFile::slice} 返回视图，而不是再次复制字节，因此每次诊断切片不需要建立新的字符串。

真实的切片函数位于 \path{E:/RComplier/Rustless/include/rx/ast/source.hpp}：
\begin{lstlisting}[language=C++]
std::string_view slice(Span span) const {
  if (span.begin > span.end || span.end > text.size())
    throw std::out_of_range("invalid source span");
  return std::string_view(text).substr(span.begin, span.end - span.begin);
}
\end{lstlisting}
范围检查使越界使用能够被发现。但视图仍然借用 \texttt{text}：销毁拥有它的 \texttt{Crate} 之后，保留这个视图是不安全的。返回视图的低复制成本与寿命约束，是同一个接口的两个方面。

当前设计还保存标识符和整数的独立拼写，构树器与最终 \texttt{Crate} 也会各自建立源码对象。因此，它并非整个过程零复制的设计。将来若处理许多文件，可以由 SourceManager 统一拥有文本，节点仅保存文件身份和范围，并在必要处采用字符串唯一化。这样的扩展首先要明确多文件身份和寿命，再讨论减少复制；单文件实现尚不能被直接描述为完整的多文件源码管理器。

\section{标识符、关键字与字面量}
观察下面的程序：
\begin{lstlisting}[language=Rustless]
fn main() {
    let mut value_1: i32 = 0x2a_i32;
    let binary = 0b1010u32;
    let octal = 0o12usize;
    let decimal = 10isize;
    let ready = true;
}
\end{lstlisting}
\texttt{value\_1} 是普通标识符。\texttt{fn}、\texttt{let} 与 \texttt{mut} 具有各自的词法类别，不应被混入普通名字。\texttt{true} 和 \texttt{false} 表示布尔字面量。整数的进制前缀、下划线分隔和后缀，则共同组成一个完整词法单元。

\path{E:/RComplier/Rustless/grammar/Lexer.g4} 中使用 ASCII 起始字符与后续字符规则；单独的 \texttt{\_} 有专门类别。由此可见，\texttt{\_} 与 \texttt{\_value} 的角色不同。生命周期名又有前置单引号，例如 \texttt{'a}。它不是字符字面量，也不是一个普通变量名字；AST 使用 \texttt{LifetimeName} 保留这一差别。

词法规则同时保留了一组暂未在子集中实现的保留关键字。例如 \texttt{match} 已被识别为保留词，却没有相应的匹配表达式 grammar。这能防止它被当成普通标识符使用，也提醒我们区分“能识别一种词”与“已经实现对应语法”。

整数识别只回答拼写是否合法。它尚未回答数值是否适合目标位宽、无后缀整数应当采用什么类型，或者一元负号怎样与整数范围共同解释。\texttt{-1} 也不是一个包含负号的整数 token：负号形成单目运算，整数载荷保存幅值的拼写。这样的划分让词法结构、语法结构和数值判断各有归属。

下面是 Lexer 中真实的整数规则摘录。它保留错误的长数字拼写，防止 \texttt{123abc} 被轻易拆成一个整数与一个标识符。
\begin{lstlisting}
INTEGER_LITERAL
    : (DECIMAL | BINARY | OCTAL | HEXADECIMAL) INTEGER_SUFFIX?
    ;
INVALID_NUMBER
    : [0-9] ASCII_CONTINUE*
    ;
fragment INTEGER_SUFFIX: 'i32' | 'u32' | 'isize' | 'usize';
\end{lstlisting}
这类规则的价值，在于错误输入仍然保有连续的范围与可解释的类别。后续诊断能够指出完整的非法数字，而不必根据偶然产生的多个 token 猜测用户原来的输入。

\section{空白、注释与非法词法单元}
空白和注释不改变这个程序的主要 AST 形状：
\begin{lstlisting}[language=Rustless]
fn main() {
    // A line comment.
    let value = 1 /* outer /* inner */ outer */ + 2;
}
\end{lstlisting}
当前 Lexer 将正常空白、单行注释和块注释放在隐藏通道中。token stream 仍能保存它们，parser 读取通常语法通道时则无需把每个空白加入产生式。AST 构造器从范围访问原文，因此即使语法结构不保存每条注释，原始文本也仍然保存在源码对象中。

嵌套块注释需要正确配对内部的开始与结束标记。当前规则通过递归的注释片段表达嵌套，且单独识别未闭合注释。它不能在遇到内部 \texttt{/*} 后，随意把这个开始标记退回成普通注释文本，否则错误的嵌套输入可能被误接受。

下列规则来自同一 Lexer 文件：
\begin{lstlisting}
WHITESPACE: SPACE_TEXT -> channel(HIDDEN);
LINE_COMMENT: LINE_COMMENT_TEXT -> channel(HIDDEN);
BLOCK_COMMENT: BLOCK_COMMENT_TEXT -> channel(HIDDEN);
UNTERMINATED_BLOCK_COMMENT: UNTERMINATED_COMMENT_TEXT;
ERROR_CHAR: .;
\end{lstlisting}
非法单元留在正常通道，正常注释进入隐藏通道。二者需要不同的安排：正常注释可以不参与语法判断，非法字符却不能静默消失。\texttt{lex\_source} 还会检查非法生命周期、非法字符字面量、非法数字、未闭合注释和未知字符类别，并产生诊断。

这并不意味着当前前端具有任意错误恢复能力。生成 parser 可以尝试恢复以继续报告错误，但 \texttt{parse\_source} 在存在诊断时停止发布 AST。消费者得到的树来自没有报告解析错误的路径，因此不必把 parser 的临时恢复节点当成可靠语法结构。

\section{FileId、Span 与半开区间}
范围采用半开区间 $[\mathrm{begin},\mathrm{end})$：包含起点，不包含终点。长度可以直接计算为终点减起点，相邻的两个范围也可以在同一终点和起点处连接。空范围自然表示为起点等于终点，而不必引入特殊负数。

实际定义位于 \path{E:/RComplier/Rustless/include/rx/diagnostic.hpp}：
\begin{lstlisting}[language=C++]
struct FileId {
  std::uint32_t value{};
  auto operator<=>(const FileId &) const = default;
};
struct Span {
  std::size_t begin{};
  std::size_t end{};
  std::size_t line{1};
  std::size_t column{};
  FileId file{};
};
\end{lstlisting}
\texttt{FileId} 用于说明范围来自哪个文件。当前单编译单元的 \texttt{SourceFile} 默认使用文件身份一；多文件分配仍需要进一步的管理器。\texttt{line} 从一开始，\texttt{column} 从零开始；打印给读者时可以再转换显示习惯。

字节偏移和字符索引不能普遍互换。ASCII 字符的 UTF-8 编码通常占一个字节，汉字等字符可能占多个字节。ANTLR 输入索引以解码后的字符为单位，AST 的切片则使用拥有的 UTF-8 字节串。\texttt{SourceFile} 为字符起点建立偏移表，并使用该表转换范围。

下图把字符编号与字节偏移放在同一段文本下方，呈现上述转换关系。这个例子只用于源码坐标教学；当前 grammar 不接受这些汉字作为语言源码中的标识符或注释文本。
\begin{figure}[htbp]
\centering
\begin{tikzpicture}[font=\small]
\node[anchor=east] at (0,0) {字符};
\node[anchor=east] at (0,-0.7) {字符起点};
\node[anchor=east] at (0,-1.4) {字节起点};
\foreach \x/\s/\c/\b in {1/中/0/0,2/文/1/3,3/+/2/6,4/1/3/7} {
  \node[draw,minimum width=14mm,minimum height=8mm] at (\x*1.5,0) {\s};
  \node at (\x*1.5,-0.7) {\c};
  \node at (\x*1.5,-1.4) {\b};
}
\end{tikzpicture}
\caption{字符位置通过偏移表映射到 UTF-8 字节位置。}
\end{figure}
转换函数的真实摘录如下，来自 \path{E:/RComplier/Rustless/include/rx/ast/source.hpp}：
\begin{lstlisting}[language=C++]
Span from_character_span(Span span) const {
  // ANTLR indexes UTF-32 characters; syntax spans index owned UTF-8 bytes.
  span.begin = characters.at(span.begin);
  span.end = characters.at(span.end);
  span.file = id;
  return span;
}
\end{lstlisting}
构树后直接使用范围切片，便不需要每次诊断重新扫描 UTF-8。代价是保存偏移表和行起点表。当前实现按输入字节建立这些表；它不是一个独立、完备的 Unicode 验证与归一化系统。扩展语言字符集时，还必须审查非法编码、字符分类与诊断列宽等需求。

\section{整体范围与局部范围}
考虑一个类型错误：
\begin{lstlisting}[language=Rustless]
fn main() {
    let value = 1 + true;
}
\end{lstlisting}
诊断可以标注整个二元表达式，说明组合不合法；也可以标注右操作数，说明需要的类型与得到的类型不同；还可以标注 \texttt{+}，说明该操作没有适合当前类型的规则。这些标注都合理，但所需范围不同。

\texttt{BinaryExpr} 因此具有 \texttt{operator\_span}，外层 \texttt{Expr} 仍具有覆盖整个表达式的 \texttt{span}。名字对象也保存自己的范围。位置不是一份泛泛的“错误行号”，而是结构中可以按需要选择的精确来源。

下图呈现整个表达式、运算符和右操作数的三种范围，说明前述诊断怎样选择不同的标注重点。
\begin{figure}[htbp]
\centering
\begin{tikzpicture}[>=Latex,font=\small]
\node[font=\ttfamily\large] at (3,0) {1 + true};
\draw[very thick] (0.8,-0.5)--(5.2,-0.5);
\node[anchor=west] at (5.5,-0.5) {整个 Expr};
\draw[thick] (2.2,-1)--(2.7,-1);
\node[anchor=west] at (5.5,-1) {运算符};
\draw[thick] (3.1,-1.5)--(5.2,-1.5);
\node[anchor=west] at (5.5,-1.5) {右操作数};
\end{tikzpicture}
\caption{一个语法结构可以提供多个有明确用途的源码范围。}
\end{figure}
保持局部范围也有维护上的意义。若以后为诊断增加高亮，运算规则不需要重新寻找 \texttt{+} 字符；若重新打印结构，整数拼写与名字范围仍可以被独立观察。源码信息应随着所描述的结构保存，而不是在最后输出错误时临时补造。

\section{生成代码与输入程序}
ANTLR 是根据语法描述生成词法分析器和语法分析器的工具。\texttt{.g4} 文件描述 token 和产生式，生成器据此建立 C++ 类；运行时库提供输入流、token stream、语法上下文和错误监听等公共机制。

项目的 \path{E:/RComplier/Rustless/CMakeLists.txt} 在构建目录中适配语法名称，再生成 \texttt{RxLexer.cpp}、\texttt{RxParser.cpp} 及相应头文件。名称适配是为了避免 C++ 目标中 \texttt{Lexer} 与 \texttt{Parser} 的保留名称冲突。源目录中的 grammar 仍保持独立描述，不需要为每个测试文件建立一套 grammar。

生成的 \texttt{RxParser} 类包含各种规则的上下文类型。运行时处理不同测试程序时，产生的是这些类的不同对象及不同树形实例。类定义由语法决定，实例内容由输入程序决定。理解这一区别，可以避免把 \texttt{generated} 中的文件误认为某个测试程序的专用 AST。

下图将生成步骤与处理程序的步骤分开，呈现前述“定义类”和“建立实例”的关系。
\begin{figure}[htbp]
\centering
\begin{tikzpicture}[>=Latex,node distance=13mm,
  box/.style={draw,rounded corners,font=\small,text width=34mm,align=center}]
\node[box] (grammar) {\texttt{.g4}\\语法描述};
\node[box,right=of grammar] (gen) {ANTLR 工具\\生成 C++ 类};
\node[box,below=of gen] (runtime) {已编译分析器\\与 runtime};
\node[box,left=of runtime] (program) {某次输入\\Rustless 程序};
\node[box,right=of runtime] (tree) {该次输入的\\parse tree};
\draw[->] (grammar)--(gen);
\draw[->] (gen)--(runtime);
\draw[->] (program)--(runtime);
\draw[->] (runtime)--(tree);
\end{tikzpicture}
\caption{语法定义决定生成组件，输入程序决定运行时建立的实例。}
\end{figure}
生成组件与手写构造器之间仍需清楚的边界。ANTLR 的上下文服务于识别过程，Rustless AST 服务于后续语言分析。下一章将观察二者怎样连接，以及构造器如何把规则层次转成语言概念。

\chapter{从语法分析树到 AST}
\label{chap:front-builder}

\section*{导读}
语法分析树记录一段输入怎样满足 grammar，AST 则记录这段程序由哪些语言结构组成。二者可以描述同一个输入，却不需要有相同的节点层次。本章从两种树的差别出发，沿着优先级、结合性、后缀链和语法上下文阅读构造器。重点是理解一个结构怎样被建立，以及哪些依赖应当留在解析边界。掌握这种阅读方法以后，读者可以在不运行整个编译器的情况下推导一段代码的 AST 形状。

\section{Parse tree 与 AST}
在 grammar 中，加法规则引用乘法规则，乘法规则引用转换规则，转换规则继续引用单目和后缀规则。这些层次很适合表达优先级，因为每一层限制自己可以包含哪些更紧密结合的表达式。但同一个整数也会通过这些中间规则向上到达 \texttt{expression}。

后续分析通常不需要逐层回答“这是 additiveExpression，还是 multiplicativeExpression 的一条通过路径”。它需要知道当前节点是整数、路径、二元运算还是调用。因此，构造器越过仅用于组织 grammar 的中间层，把真正有意义的结构保存为 AST 载荷。

下图对照 \texttt{1 + 2} 的部分语法规则树与最终 AST，呈现前文所述的层次压缩。为便于阅读，语法规则树只画出与论点有关的层级。
\begin{figure}[htbp]
\centering
\begin{tikzpicture}[>=Latex,node distance=8mm,
  box/.style={draw,rounded corners,font=\small,align=center}]
\node[box] (pe) {expression};
\node[box,below=of pe] (pa) {additiveExpression};
\node[box,below left=10mm and 2mm of pa] (p1) {multiplicativeExpression\\literal 1};
\node[box,below right=10mm and 2mm of pa] (p2) {multiplicativeExpression\\literal 2};
\draw[->] (pe)--(pa); \draw[->] (pa)--(p1); \draw[->] (pa)--(p2);
\node[box,right=44mm of pa] (ae) {Expr\\BinaryExpr: Add};
\node[box,below left=10mm and 0mm of ae] (a1) {Expr\\Integer 1};
\node[box,below right=10mm and 0mm of ae] (a2) {Expr\\Integer 2};
\draw[->] (ae)--(a1); \draw[->] (ae)--(a2);
\end{tikzpicture}
\caption{规则树与 AST 使用不同层次服务不同任务。}
\end{figure}
抽象并不意味着删除一切与拼写有关的信息。当前 AST 保留 \texttt{ParenExpr} 与 \texttt{ParenType}，也保存整数原拼写。这些信息可以服务源码观察和精确诊断。需要忽略括号的处理过程通过专门接口展开；它不必迫使其他使用者失去括号的存在信息。

这种分离还有一项重要结果：AST 的公开定义不包含 ANTLR 上下文类型。若调整 grammar 的表达式分层，许多变化只发生在构造器；二元表达式的左右字段及其使用者可以保持稳定。稳定的基础应当是语言结构，而不是生成器当前产生的某个类名。

\section{优先级如何形成父子关系}
\begin{lstlisting}[language=Rustless]
fn main() {
    let first = 1 + 2 * 3;
    let second = (1 + 2) * 3;
    let test = 1 + 2 < 4 && true;
}
\end{lstlisting}
乘法比加法结合得紧密，因此第一条初始化的根是加法，右子节点是乘法。第二条初始化的根是乘法，其左子节点保留括号，括号中再包含加法。第三条初始化先形成整数计算和比较，再把比较结果作为逻辑运算的输入。

\path{E:/RComplier/Rustless/grammar/Parser.g4} 的以下真实产生式说明了这项层次关系：
\begin{lstlisting}
additiveExpression
    : multiplicativeExpression (additiveOperator multiplicativeExpression)*
    ;
multiplicativeExpression
    : castExpression (multiplicativeOperator castExpression)*
    ;
castExpression
    : unaryExpression (AS typeRef)*
    ;
\end{lstlisting}
加法的每个操作数首先是一个完整的乘法表达式，而乘法的每个操作数又首先是一个完整的转换表达式。构造器按该结构读取子上下文，便不必另写一份优先级数值表来重新分析相同输入。

从紧密到松散，当前表达式 grammar 依次组织后缀、单目、转换、乘除余数、加减、移位、按位与、按位异或、按位或、比较、逻辑与、逻辑或和赋值。这个顺序是语言规则，不是 Visitor 重载的先后顺序。即使 C++ 处理函数的排列发生变化，只要 AST 已正确建立，优先级也不会随之改变。

这样的设计使一个重要决定只有一个主要来源。grammar 确定结合关系，构造器保存它，后续过程读取它。若后端再用运算符文本重新判断优先级，就会制造两份需要同时维护的解释。

\section{结合性与左折叠}
优先级说明不同运算之间怎样嵌套，结合性则说明同一层重复运算怎样组成树。对 \texttt{10 - 3 - 2}，左结合产生 \texttt{(10 - 3) - 2}；若错误地右结合，结果结构会变为 \texttt{10 - (3 - 2)}，语义也随之改变。

构造器用一个逐步累积的结果实现左折叠。它先接收第一个操作数，然后每遇到下一个操作数，就把已经建立的结果作为新节点的左孩子。以下是 \path{E:/RComplier/Rustless/src/frontend.cpp} 中 \texttt{fold\_binary} 的真实删节摘录；省去了 token 与规则上下文识别，只保留建立关系的部分。
\begin{lstlisting}[language=C++]
auto operand = build_expression(part);
if (!result) {
  result.emplace(std::move(operand));
  continue;
}
auto span = merge((*result)->span, operand->span);
auto binary = make_expr(
    context,
    ast::BinaryExpr{ast::binary_operator(operator_text), operator_span,
                    std::move(*result), std::move(operand)});
binary->span = span;
result = std::move(binary);
\end{lstlisting}
\texttt{std::move} 在这里转移已经建立的子树拥有关系。它没有深复制整个左子树；新二元节点接管旧结果，随后累积变量接管新节点。这样，构树步骤与最后的树形一致，且不会因反复复制已有前缀而增加不必要的工作。

下图逐步展示左折叠怎样把三个操作数组成两层二元节点，对应上述累积过程。
\begin{figure}[htbp]
\centering
\begin{tikzpicture}[>=Latex,node distance=10mm,
 box/.style={draw,rounded corners,font=\small,align=center,text width=30mm}]
\node[box] (one) {第一步\\\texttt{10}};
\node[box,right=of one] (two) {第二步\\\texttt{Sub(10, 3)}};
\node[box,right=of two] (three) {第三步\\\texttt{Sub(previous, 2)}};
\draw[->] (one)--(two); \draw[->] (two)--(three);
\end{tikzpicture}
\caption{旧结果成为下一层二元表达式的左子树。}
\end{figure}
赋值 grammar 采用另一种结构：右侧递归到整个表达式，所以 \texttt{a = b = 1} 的语法树按右侧嵌套。本例仅用于说明构树，不能据此断言赋值链通过语义检查。转换链 \texttt{x as u32 as i32} 则逐步包裹前一个结果：先形成转换到 \texttt{u32} 的节点，再以它为输入形成转换到 \texttt{i32} 的节点。

这几种构造方式不应被统一成没有说明的“不断添加操作数”。每一步到底是在左侧累积、递归右侧，还是包裹先前结果，是语言结构必须保留的差别。

\section{后缀链怎样包裹已有表达式}
下例使用具有实际声明的字段与方法：
\begin{lstlisting}[language=Rustless]
struct Pair { values: [i32; 2] }
impl Pair {
    fn get(&self) -> i32 { self.values[0] }
}
fn main() {
    let p = Pair { values: [1, 2] };
    let a = p.values[0];
    let b = p.get();
}
\end{lstlisting}
\texttt{p.values[0]} 的起点是路径 \texttt{p}。读取点号后，构造器建立字段节点，把 \texttt{p} 放入 \texttt{base}。读取方括号后，又建立索引节点，把整个字段节点放入新的 \texttt{base}。最终根节点是 IndexExpr，而不是一条平铺的“路径及附件”记录。

方法调用 \texttt{p.get()} 则把 \texttt{p} 作为 \texttt{receiver}，将方法名字与普通参数分开保存。普通调用的起点是 \texttt{callee}；它可以本身是一个表达式，例如局部变量所持有的函数项。相似的表面标点因此并不意味着所有后缀都应使用同一个字段集合。

下面是 \texttt{apply\_call} 的真实摘录，来自 \path{E:/RComplier/Rustless/src/frontend.cpp}：
\begin{lstlisting}[language=C++]
auto span = merge(base->span, source_span(context));
auto call = make_expr(
    context, ast::CallExpr{std::move(base), call_arguments(context)});
call->span = span;
return call;
\end{lstlisting}
范围合并说明完整调用包括它的被调用表达式，而不只包括参数括号。子树移动说明调用接管已有 base 的所有权。构造器的后缀循环按源顺序应用这些变换，每一步都产生下一步可以接收的完整表达式。

下图呈现调用、字段与索引依次包裹已有 base 的过程。图中的复杂链只是语法结构示意，不假定 \texttt{f} 与它返回的类型已经声明。
\begin{figure}[htbp]
\centering
\begin{tikzpicture}[>=Latex,node distance=8mm,
 box/.style={draw,rounded corners,font=\small,align=center,text width=27mm}]
\node[box] (p) {\texttt{f}\\PathExpr};
\node[box,right=of p] (c) {\texttt{f(1)}\\CallExpr};
\node[box,right=of c] (f) {\texttt{f(1).items}\\FieldExpr};
\node[box,right=of f] (i) {\texttt{f(1).items[0]}\\IndexExpr};
\draw[->] (p)--(c); \draw[->] (c)--(f); \draw[->] (f)--(i);
\end{tikzpicture}
\caption{每一步后缀都接收前一步形成的表达式。}
\end{figure}

\section{块、分号和尾表达式}
块包含语句序列，也可以有一个没有分号的尾表达式。
\begin{lstlisting}[language=Rustless]
fn first() -> i32 {
    let x = 1;
    x + 2
}
fn second() {
    let x = 1;
    x + 2;
}
\end{lstlisting}
第一段中的加法属于 \texttt{Block.tail}，第二段则属于语句序列中的 ExprStmt。若只在某个通用标志中记录“块最后没有分号”，每个使用者仍要重新找到最后节点并解释这个标志。当前构造器把该关系直接变成不同的结构位置。

对于 grammar 已形成的语句，构造器逐项建立 Stmt；对于明确的尾表达式，构造器建立可选 tail。某些无分号的块类表达式先经语句规则识别，随后在构造阶段从最后一条语句移入 tail。真实的相关代码如下：
\begin{lstlisting}[language=C++]
if (context->statementExpression()) {
  block->tail.emplace(build_expression(context->statementExpression()));
} else if (!block->statements.empty()) {
  auto *last = std::get_if<ast::ExprStmt>(&block->statements.back()->data);
  if (last && last->termination == ast::Termination::Omitted) {
    block->tail.emplace(std::move(last->expression));
    block->statements.pop_back();
  }
}
\end{lstlisting}
这是 \texttt{build\_block} 的删节摘录，来源仍是 \path{E:/RComplier/Rustless/src/frontend.cpp}。它解释的是 AST 的归属整理，尚未决定整个块的类型。块值与函数返回值是否匹配，需要后面的语义判断。

\section{类型与词法歧义的构造边界}
某些标点在不同上下文中有不同用途。\texttt{Vec::<Vec<i32>>::new()} 含有两层泛型闭合符；\texttt{x >> 1} 含有一个右移运算；\texttt{x >>= 1} 又是复合赋值。把所有相邻的 \texttt{>} 随意拼成右移，会破坏嵌套类型；无条件拆分右移也会使运算规则难以识别原来的整体 token。

当前 Lexer 在读取 \texttt{>} 后进入专门模式，用不同 token 类型保留一个最大标点组合中的后续部分。parser 的 \texttt{genericClose} 读取单个闭合部分，\texttt{shiftRight} 与赋值规则则读取相应完整组合。来自两个不同 token 的部分不能仅凭相邻位置重新组成一个运算符。这项细节被限制在 Lexer 与 grammar 中；AST 保存的仍然是 \texttt{BinaryOp::Shr} 或 \texttt{CompoundAssignOp::Shr}。

引用类型中的 \texttt{&&} 也需要展开为两层，而不是一种新的“二重引用”节点。对语法片段 \texttt{\&\&'a mut i32}，外层是不可变引用，生命周期和可变性属于内层。构造器首先建立内层，再用外层包裹它，并调整内层的源码起点。这样，嵌套引用直接复用 ReferenceType，不需要让语义阶段再解释一个特殊旗标。

转换后的 \texttt{<} 同样可能被看作类型参数的开始。写为 \texttt{x as (i32) < y} 时，括号明确界定转换类型；grammar 中的 closed 类型和表达式规则负责表达这类边界。教读者理解这些规则，比让后续阶段看到“含糊的转换字符串”后再猜测意图更清晰。

\section{普通、条件与语句表达式}
条件中的 \texttt{if value \{ ... \}} 需要区分条件路径与随后分支块。若无界定地允许 \texttt{value \{ ... \}} 在该处成为结构体构造，parser 就可能把分支大括号吞进条件。当前 condition 规则因此限制未加界定的结构体构造，而括号、数组、调用参数等界定结构中的表达式可以使用普通规则。

条件中的 break 操作数还有相似问题：若它以一个裸块开始，块可能与正在寻找的分支边界混淆。因此 grammar 对该起点作出限制，同时保持后续操作数的正常结合。这些规则是为消除具体语法歧义而设，不应被推广成“条件表达式的语义总是不同”。

语句开头也有自己的限制。一个块类表达式可以直接组成语句；若它还要继续形成后缀表达式，当前 grammar 要求先经过点号后缀。用括号界定的表达式则有更清楚的起点。这些选择说明了为什么 Parser.g4 同时存在普通、条件和语句表达式家族。

构造器根据规则名称识别这些家族，却把相同的实际结构转成同一套 AST 载荷。条件中的加法、语句中的加法与普通加法都成为 BinaryExpr。于是语义规则只需要处理“加法表达式”，不必复制三套以 parser 上下文命名的类型。

这一点体现了抽象边界的作用：grammar 可以因识别需求更细分，AST 则按已确定的语言形状汇合。若让三个语法上下文变成三个语义节点，每当修改整数运算规则，就必须重复修改原本相同的判断。

\section{构树成功与错误处理}
\texttt{AstBuilder} 只在 parser 没有报告错误时建立并发布树。节点身份和源码范围由 \texttt{stamp} 统一填写，表达式的具体载荷则由 \texttt{make\_expr} 装入 variant。真实摘录来自 \path{E:/RComplier/Rustless/src/frontend.cpp}：
\begin{lstlisting}[language=C++]
template <typename Node> void stamp(Node &node, ParseContext *context) {
  node.id = ast::NodeId{next_id_++};
  node.span = source_span(context);
}
template <class Payload>
ast::ExprPtr make_expr(ParseContext *context, Payload payload) {
  auto expression = ast::box<ast::Expr>(ast::Expr::Data{std::move(payload)});
  stamp(*expression, context);
  return expression;
}
\end{lstlisting}
集中处理共同信息能够减少遗漏，但它不代替各载荷的构造责任。二元载荷仍需要明确提供两个子节点，函数仍需要提供函数体。构造器遇到已经假定存在却找不到的操作数时，抛出内部逻辑错误；这种错误说明手写构造逻辑或 grammar 适配出了问题，与用户程序的正常语法诊断不同。

生成 parser 的类在输入流、token stream 和 parser 对象的寿命范围内使用。构树器把语言所需的数据移入自己的 AST，\texttt{parse\_source} 最后返回拥有树的结果。后续使用者不需要保存 parser 来读取表达式。下一章将从这些构造要求出发，研究 AST 的类型和拥有关系。

\chapter{节点身份、具体载荷与所有权}
\label{chap:front-ownership}

\section*{导读}
二元表达式需要左右两个操作数，返回表达式可以没有返回值，函数调用则需要一个有序参数序列。让这些差别停留在注释和约定中，处理过程就必须反复检查它们。本章讨论如何把结构要求写入 C++ 类型，并解释节点身份、活动载荷和子节点拥有关系各自承担的职责。这里的类型安全首先指语法结构的可表达性；表达式的语言类型仍由语义阶段判断。

\section{相同拼写与不同出现}
\begin{lstlisting}[language=Rustless]
fn main() {
    let first = 1 + 2;
    let second = 1 + 2;
}
\end{lstlisting}
两次加法有相同运算符和相同字面量，却位于不同源码范围中，也属于不同变量的初始化。若一条语义事实只用字符串 \texttt{1 + 2} 作为键，便无法确定它属于哪一次出现。\texttt{NodeId} 为节点提供身份，使位置相同与结构相同不必被误认为同一节点。

项目使用带标签的整数包装，而非把每种身份都声明为普通 \texttt{uint32\_t}。来自 \path{E:/RComplier/Rustless/include/rx/ast/node.hpp} 的摘录如下：
\begin{lstlisting}[language=C++]
template <class Tag> struct Id {
  std::uint32_t value{};
  auto operator<=>(const Id &) const = default;
};
struct NodeTag;
using NodeId = Id<NodeTag>;
\end{lstlisting}
标签参与 C++ 类型，因此节点身份与使用另一个标签的声明身份不能无意互换。它们可能使用相同的数值表示，却不是相同的概念。源码范围也不适合充当唯一身份：不同层次节点的范围可以重叠，括号包装与内部表达式还可能具有紧密相邻的范围。

当前身份在一次构树过程中分配，零通常表示尚未设置。它不是跨编译持久编号，也不保证按源位置排序；某些外层节点是在子结构完成后才分配。若未来支持增量分析，需要另行设计版本与持久关联，不能仅凭当前递增整数推导出这些能力。

\section{Node 与承载节点}
\texttt{Node} 保存各类有身份节点共同需要的信息：
\begin{lstlisting}[language=C++]
struct Node {
  NodeId id{};
  Span span{};
};
\end{lstlisting}
\texttt{Expr : Node} 继承这两项信息，但一个 Expr 的具体形状由其 \texttt{Data} 决定。我们把有身份的外层称为承载节点，即 owner，把描述形状的数据称为具体载荷，即 payload。例如，BinaryExpr 保存运算和左右孩子；它本身没有继承 Expr。

下图把表达式身份与范围放在外层，把具体语法字段放在活动载荷中，呈现这两层共同参与访问的关系。
\begin{figure}[htbp]
\centering
\begin{tikzpicture}[>=Latex,node distance=9mm,font=\small,
 box/.style={draw,rounded corners,align=center}]
\node[box,text width=76mm] (owner) {Expr：NodeId、Span 与 Data};
\node[box,text width=76mm,below=of owner] (payload) {活动载荷 BinaryExpr\\op、operator span、left、right};
\node[box,below left=10mm and 3mm of payload] (left) {左 Expr\\自身身份与载荷};
\node[box,below right=10mm and 3mm of payload] (right) {右 Expr\\自身身份与载荷};
\draw[->] (owner)--node[right]{拥有载荷}(payload);
\draw[->] (payload)--node[left]{left}(left);
\draw[->] (payload)--node[right]{right}(right);
\end{tikzpicture}
\caption{身份层与具体形状层分别表达不同的结构信息。}
\end{figure}
不是每种节点都需要两种 C++ 结构体。Function 的字段直接描述一种确定形状，因此 Function 自身既提供身份，也提供具体内容。统一访问接口仍可把它同时作为 owner 与 payload 传入。对读者而言，关键在于理解接口中的职责，而不是要求所有节点机械地多包一层。

\section{从通用字段集合到 variant}
早期 AST 常采用一个枚举标签和一组通用字段。下面是这种方式的教学缩略示意，不是当前项目定义：
\begin{lstlisting}[language=C++]
struct OldExpr {
  Kind kind;
  std::string text;
  std::vector<ExprPtr> operands;
  TypePtr cast_type;
  bool flag;
};
\end{lstlisting}
这种结构允许表示整数，却同时携带转换类型；也允许把 kind 设为 Binary 后仍不给两个操作数。\texttt{text} 与 \texttt{flag} 的含义随种类变化，读者必须查找外部约定才能理解字段。每个消费者又都需要自行确认这些约定。

具体载荷将关系改写为命名字段。真实 BinaryExpr 定义如下，来自 \path{E:/RComplier/Rustless/include/rx/ast/expr.hpp}：
\begin{lstlisting}[language=C++]
struct BinaryExpr {
  BinaryOp op;
  Span operator_span;
  ExprPtr left;
  ExprPtr right;
};
\end{lstlisting}
二元表达式不会额外拥有 callee 或字段初始化列表。Expr 的 variant 在同一时刻保存一种载荷，并管理相应对象的构造与析构。kind 与数据形状不再是两份手工维护的声明：当前活动的 C++ 类型就说明形状。

这项约束并不检查 \texttt{left} 的语言类型是否支持 \texttt{op}。\texttt{1 + true} 依然可以建立 BinaryExpr，因为它确实具有两个操作数和一个加法运算符。结构安全使语义检查面对清楚的输入，语义检查再判断这些输入能否组合。让两层承担各自能够表达的关系，比用一个“合法”标志概括全部判断更容易维护。

\section{必需、可选与有序子结构}
一个二元表达式缺少右操作数，通常没有合理的完整 AST 含义。一个 \texttt{return} 没有返回值，却可能是合法的写法。一个调用没有参数，还应有被调用表达式。这几种情况应当分别表示。

必需孩子使用 \texttt{Box<T>}；可选孩子使用 \texttt{optional<Box<T>>}；有序序列使用 \texttt{vector<Box<T>>}。下面的真实载荷定义显示区别：
\begin{lstlisting}[language=C++]
struct CallExpr {
  ExprPtr callee;
  std::vector<ExprPtr> arguments;
};
struct ReturnExpr {
  std::optional<ExprPtr> value;
};
\end{lstlisting}
\begin{lstlisting}[language=Rustless]
fn zero() {}
fn add(a: i32, b: i32) -> i32 { a + b }
fn main() {
    zero();
    let result = add(1, 2);
    return;
}
\end{lstlisting}
空参数序列是调用的一个普通状态，空 callee 却不是。没有返回值由 optional 的缺席表示；\texttt{return ()} 则有一个实际 UnitExpr。二者可以在语义上产生相关结果，但在语法结构上仍有区别。

下图用三种边呈现必需、可选与序列关系，对应上述 C++ 容器，而不是程序的运行时对象引用。
\begin{figure}[htbp]
\centering
\begin{tikzpicture}[>=Latex,node distance=12mm,font=\small,
 box/.style={draw,rounded corners,align=center}]
\node[box] (c) {CallExpr};
\node[box,below left=13mm and 10mm of c] (cal) {callee：必需 Expr};
\node[box,below right=13mm and 10mm of c] (args) {arguments：有序 Expr 序列};
\node[box,right=34mm of c] (r) {ReturnExpr};
\node[box,below=of r] (rv) {value：可选 Expr};
\draw[->,thick] (c)--(cal); \draw[->,double] (c)--(args);
\draw[->,dashed] (r)--(rv);
\end{tikzpicture}
\caption{子结构是否存在以及存在多少项，由其类型直接表达。}
\end{figure}
这里的 C++ Box 用于拥有 AST 子节点，与源语言 \texttt{Box<i32>} 的运行时容器是不同概念。前者解决编译器数据结构的寿命，后者需要语言语义、内存布局和运行时操作。相同名称不能替代对用途的说明。

\section{递归类型与移动所有权}
Expr 的载荷可以拥有其他 Expr；TypeSyntax 的数组长度又可以拥有 Expr。这形成递归定义。直接把所有孩子按值嵌入对象会导致无法确定有限的对象大小，因此孩子需要间接存储。

当前 Box 使用独占指针，禁止复制，并在构造时拒绝空指针。以下为 \path{E:/RComplier/Rustless/include/rx/ast/node.hpp} 的真实删节摘录：
\begin{lstlisting}[language=C++]
template <class T> class Box {
public:
  explicit Box(std::unique_ptr<T> value)
      : value_(value.release(), [](T *pointer) { delete pointer; }) {
    if (!value_)
      throw std::invalid_argument("a required AST child cannot be null");
  }
  Box(Box &&) noexcept = default;
  Box &operator=(Box &&) noexcept = default;
  Box(const Box &) = delete;
  Box &operator=(const Box &) = delete;
  const T *operator->() const { return value_.get(); }
private:
  std::unique_ptr<T, void (*)(T *)> value_;
};
\end{lstlisting}
独占拥有使树的销毁路径明确：父节点销毁时，其孩子也沿拥有边销毁。构造时移动子树，则把责任交给新的父节点。它不会使两个父节点同时负责删除同一个孩子。若需要跨树共享结构，必须另行选择相应的拥有模型，而不能随意复制 Box。

递归头文件还存在不完整类型问题。某个头文件可能只知道 \texttt{struct Block;}，却需要保存 \texttt{Box<Block>}。当前删除器在 T 已完整的位置创建，并以函数指针保存，销毁持有者时不要求每个递归头文件都包含全部节点定义。这帮助将语法族分开，但通常增加一个删除器指针和间接删除调用的成本。

构造时非空并不意味着任意时刻的任意 Box 都可解引用。移动后原对象可能为空；它应被销毁或重新赋值，不应再被当作有效孩子读取。完整 AST 的必需孩子应处在有效拥有状态。理解 moved-from 对象的限制，才能准确说明这项类型约束。

\section{只读发布与可构造性}
前端发布 \texttt{unique\_ptr<const Crate>}，正常读取路径因此获得 const 树。Box 的 const 解引用返回 const 子节点，vector 中的载荷与子结构也通过 const 路径观察。Visitor 接口进一步只接收 const owner 和 const payload。

只读发布让多种分析过程可以共享结构，而不必担心某个打印步骤顺手改写运算符或移动孩子。语义事实独立保存，也便于描述“读树并建立结论”这一过程。若某个后续变换需要改写程序，宜为它明确选择修改接口或新的表示，而不是让普通查询隐含修改能力。

不过，节点字段仍可在手工构造时被赋值，id 和 span 也不是构造参数中的强制条件。类型系统已经排除了一部分非法形状，构树器和一致性检查仍负责其他关系，例如身份唯一、范围有效和来源归属。这种分层约束应当如实说明；不能因为接口只读，就推导出每一个手工构造的对象都已经正确。

\section{存储方案的取舍}
variant 对象需要足够的空间容纳最大的活动载荷，并保存判别信息。它通常不为自己的活动载荷单独申请堆空间，但载荷中的 string、vector 以及 Box 孩子仍可能分配。因而“使用 variant”与“所有节点紧凑连续”之间没有必然关系。

逐节点 Box 分配容易表达树拥有关系，却带来分配次数和间接访问。arena 可以把一批节点放入集中管理的存储区域，用整体寿命减少个别删除操作；索引式节点可以把边表示为编号，并使用数组管理对象。但这些方式还需要回答地址稳定、不同载荷尺寸、容器内部分配和对象析构的问题。

下图比较逐节点拥有与集中分配的访问路径，用来解释上述分配与局部性取舍。集中区域只是后续设计方向，不是当前 AST 的实际布局。
\begin{figure}[htbp]
\centering
\begin{tikzpicture}[>=Latex,font=\small,
 box/.style={draw,rounded corners,minimum width=22mm,minimum height=8mm}]
\node at (2,1.4) {当前：独立拥有};
\node[box] (p) at (2,0.6) {父节点};
\node[box] (a) at (0.7,-0.6) {孩子 A};
\node[box] (b) at (3.4,-1.1) {孩子 B};
\draw[->] (p)--(a); \draw[->] (p)--(b);
\node at (8,1.4) {可能方向：集中存储};
\node[draw,minimum width=58mm,minimum height=22mm] at (8,-0.2) {};
\node[box] at (6.5,-0.2) {父};
\node[box] at (8,-0.2) {A};
\node[box] at (9.5,-0.2) {B};
\end{tikzpicture}
\caption{不同拥有方式改变分配与访问路径，需要结合寿命和测量选择。}
\end{figure}
更紧凑的布局可能有利于遍历，也可能使构造和改写更复杂。虚函数对象也有各自的存储与分派代价。应根据实际节点尺寸、分配次数、访问热点与维护需求比较，而不应把任何一种机制当作普遍最优。

本章明确了所有节点共同遵守的原则。下一章将阅读名字、类型语法和声明，观察这些原则怎样落到完整语言子集中的具体字段。
